%% EAAI V3 — Engineering Applications of Artificial Intelligence
%% Document class: elsarticle (Elsevier); compiled with pdflatex.
%% At submission use option [review] for double-blind; switch to [preprint]
%% for the accepted-manuscript deposit.
\documentclass[review,12pt]{elsarticle}

\usepackage{amsmath,amssymb}
\usepackage{graphicx}
\graphicspath{{../figures/}}
\usepackage{booktabs}
\usepackage{array}
\usepackage{tabularx}
\usepackage{multirow}
\usepackage{algorithm}
\usepackage{algpseudocode}
\usepackage{siunitx}
\usepackage{xcolor}
\usepackage[hidelinks]{hyperref}

\newcommand{\bp}{\ensuremath{\,\mathrm{bp}}}
\newcommand{\todo}[1]{\textcolor{red}{[#1]}}

%% Journal information (filled at submission)
\journal{Engineering Applications of Artificial Intelligence}

%% ---------------------------------------------------------------
\begin{document}

\begin{frontmatter}

\title{A Physics-Informed Neural Network for Option Pricing under Rough
Volatility: Architecture, Dimension Scaling, and Fast Parametric Calibration
via the Markovian-Lifted PDE}

\author[cu]{Aviral Sharma\corref{cor}}
\ead{aviral98765@gmail.com}
\author[sgt]{Kamal Preet Singh}
\author[ind]{Vipul Sharma}

\cortext[cor]{Corresponding author}
\affiliation[cu]{organization={Chandigarh University},
              city={Chandigarh}, country={India}}
\affiliation[sgt]{organization={SGT University},
              city={Gurugram}, country={India}}
\affiliation[ind]{organization={Independent Researcher}, country={India}}

%% ----- Highlights (3-5 bullets, max 85 characters each) --------
\begin{highlights}
\item A PINN solves the $(n{+}1)$-dimensional Markovian-lifted rough-Heston PDE
      without a computational grid.
\item Solve-error grows sub-linearly in the lifted dimension ($<11\,\mathrm{px\,bp}$
      to $n{=}32$) while grid methods become infeasible.
\item Explicit error decomposition separates network accuracy from model
      approximation, yielding a principled rule for choosing $n$.
\item A parametric PINN over five model parameters warm-starts calibration at
      ${\sim}9\times$ fewer true-model evaluations.
\item Four stabilising design choices are identified, implemented, and ablated.
\end{highlights}

%% ----- Abstract ------------------------------------------------
\begin{abstract}
Rough-volatility models capture empirical features of equity and cryptocurrency
option markets that classical models cannot, but they carry a severe computational
penalty: the variance is non-Markovian, and there is no finite-dimensional pricing
partial differential equation (PDE) to solve.  The Markovian lift recovers a PDE
by approximating the fractional variance kernel with $n$ exponentials, at the cost
of turning a two-dimensional problem into an $(n{+}1)$-dimensional one.  We solve
this high-dimensional lifted PDE with a physics-informed neural network (PINN)---a
fully connected, $\tanh$-activated multilayer perceptron trained by minimising the
strong-form PDE residual at randomly sampled collocation points---and ask a sharp
engineering question: does the network's accuracy degrade as the lifted dimension
grows, in the same way a finite-difference grid would?  It does not.  Measuring
the network's solve-error against an independent lifted Monte Carlo at the same $n$,
and separating it cleanly from the model-approximation lift-error, we find the
solve-error grows only sub-linearly: from $5.4\pm1.0$ to $10.7\pm2.0$ price basis
points as $n$ climbs from 4 to 32, while the naive grid requires $50^{n+1}$ points
and is infeasible by $n{=}8$.  We detail four architecture and training decisions
that are essential for stable convergence: a hard-constrained initial-condition
ansatz, an integrated-variance baseline anchor, self-adaptive loss weighting, and a
causal $\tau$-curriculum.  We also build a parametric PINN that accepts the five
rough-Heston parameters as additional network inputs; while too coarse to calibrate
independently, it provides a warm start that, combined with a short local
true-model refinement, recovers full calibration quality at approximately nine times
fewer exact-pricer evaluations.  All experiments use live Bitcoin option data from
the Deribit exchange, returning a genuinely rough Hurst exponent ($H\approx0.09$)
with strong negative correlation ($\rho\approx-0.79$).
\end{abstract}

\begin{keyword}
Physics-informed neural networks \sep rough volatility \sep Markovian lift \sep
option pricing \sep curse of dimensionality \sep model calibration \sep
deep learning \sep scientific machine learning
\end{keyword}

\end{frontmatter}

% ================================================================
\section{Introduction}
\label{sec:intro}

Every derivative-pricing problem is, at its core, a PDE.  For the classical models
of quantitative finance that PDE is low-dimensional and a finite-difference grid
dispatches it in milliseconds.  Rough volatility breaks this comfortable picture.
Motivated by the empirical roughness of realised variance
\citep{gatheral2018volatility,bennedsen2022decoupling} and the steep short-maturity
skew that classical stochastic-volatility models such as \citet{heston1993closed}
cannot reproduce, and with microstructural foundations in high-frequency trading
dynamics \citep{eleuch2018microstructural}, rough-volatility models drive the
variance with a fractional noise whose Hurst exponent $H<1/2$.  That single
fact---whose lineage runs back to long-memory volatility \citep{comte1998long}---makes
the variance process neither Markovian nor a semimartingale, leaving
\emph{no} finite-dimensional pricing PDE to discretise at all
\citep{gatheral2018volatility,eleuch2019characteristic}.

The standard engineering escape is the \emph{Markovian lift}: approximate the
fractional kernel by a sum of $n$ exponentials and recover a genuine
$(n{+}1)$-dimensional PDE \citep{abijaber2019multifactor,abijaber2019lifting}.  The
lift is an approximation whose quality improves with $n$---the relative $L^2$ kernel
error falls from $25.6\%$ at $n{=}5$ to $0.39\%$ at $n{=}100$---but it couples the
approximation-order knob $n$ directly to the PDE dimension.  That coupling makes
the lifted PDE a clean engineering stress-test for mesh-free neural solvers: as we
turn $n$ up to improve model fidelity, does the solver hold?

Physics-informed neural networks (PINNs) are built for this setting.  Rather than
laying down a grid, a PINN represents the price as a neural network and trains it
to satisfy the PDE at randomly sampled collocation points
\citep{raissi2019physics}.  Because nothing in the method scales with a
grid, it sidesteps the exponential blow-up that kills finite differences in high
dimensions, in the same way as the deep Galerkin and deep BSDE solvers
\citep{sirignano2018dgm,han2018solving}.  For pricing it brings two further
engineering advantages: the physics is known exactly, so no labelled prices are
needed during training; and the trained network is differentiable in its inputs,
which we exploit to calibrate by gradient descent.

We develop, validate, and characterise a PINN for the lifted rough-Heston PDE and
test it throughout on Bitcoin (BTC) options---a liquid, high-skew market that, as
we show, calibrates to a firmly rough Hurst exponent.  Our concrete engineering
contributions are as follows.

\begin{enumerate}

\item \textbf{A complete, reproducible PINN implementation for the lifted PDE.}
We describe the architecture, training stabilisers, and collocation strategy in
full detail (Section~\ref{sec:arch}), identify the four design choices that are
individually necessary for stable convergence, and make the full PyTorch
implementation available.

\item \textbf{A dimension-scaling result.}  The network's solve-error---how well
it actually solves the PDE it is handed---grows only sub-linearly as the lifted
dimension climbs, staying below roughly eleven price basis points out to $n{=}32$,
while a naive full-tensor finite-difference grid ($50^{n+1}$ points) is already
infeasible by $n{=}8$.  We are deliberately fair about the comparison: $50^{n+1}$
is the naive grid; structured schemes do better but still degrade with $n$.  The
honest engineering verdict is that a competitive grid solve at $n{=}16$--$32$ is
\emph{impractical}, not impossible.

\item \textbf{An error decomposition and a practical rule for choosing $n$.}
We split the total pricing error into the part the network controls
(solve-error, measured against a lifted Monte Carlo at the same $n$) and the
part only more factors can fix (lift-error, measured against the exact Fourier
pricer).  Their crossover gives a concrete, maturity-dependent rule for the factor
count.

\item \textbf{A real-market calibration.}  Fit to live Deribit BTC surface data,
rough Heston returns a rough Hurst exponent and strong negative correlation.

\item \textbf{Fast recalibration, honestly stated.}  A parametric PINN plus a
warm-start hybrid recovers from-scratch calibration quality at approximately nine
times fewer exact-model evaluations; we are explicit about where the standalone
parametric network falls short.

\end{enumerate}

We are explicit about what is \emph{not} new.  \citet{papapantoleon2025timestepping}
solved the lifted rough-Heston PDE with a mesh-free network before us.  Our
contribution is a residual PINN formulation, a quantitative scaling-and-decomposition
analysis, a real-market calibration, and the recalibration hybrid.

% ================================================================
\section{Related work}
\label{sec:related}

\paragraph{Closest prior work.}
\citet{papapantoleon2025timestepping} price European options in rough-diffusion
models by recasting the pricing PDE as an energy-minimisation problem and
approximating it in a time-stepping deep gradient flow, demonstrated on the lifted
Heston model.  Their method is the closest to ours in both target and spirit;
it differs in mechanism (time-stepping variational scheme versus residual PINN)
and does not study dimension scaling, error decomposition, or the recalibration
workflow.

\paragraph{Neural solvers for rough volatility.}
\citet{jacquier2023random} cast the path-dependent pricing equations as backward
SDEs and solve them with a random-feature network, which reduces training to a
least-squares regression.  They share the target but not the mechanism, and
dimension scaling and calibration speed are not studied.

\paragraph{Deep calibration surrogates.}
The dominant approach for fast rough-volatility calibration is the supervised
surrogate of \citet{bayer2019deep} and \citet{horvath2021deep}: a network trained
offline on precomputed prices or implied volatilities, used as a fast pricing map
inside the calibration loop.  It is purely data-driven; the network never sees the
governing equation, carries no PDE consistency, and cannot warm-start a
physics-based polish the way ours can.

\paragraph{PINN training pathologies and remedies.}
PINNs are well-known to fail when loss terms sit on different scales or when
temporal causality is violated \citep{krishnapriyan2021characterizing,wang2022when}.
The two stabilisers we use---self-adaptive homoscedastic-uncertainty weights
\citep{kendall2018multitask,mcclenny2023self} and a causal $\tau$-curriculum
\citep{wang2024respecting}---address these failure modes.

\paragraph{Foundational methods.}
We build on PINNs \citep{raissi2019physics} and deep PDE solvers
\citep{sirignano2018dgm,han2018solving}; on rough volatility
\citep{gatheral2018volatility,eleuch2019characteristic} and its multifactor
Markovian lift \citep{abijaber2019multifactor,abijaber2019lifting}; and on
classical transform pricing \citep{heston1993closed,carr1999option} for our
independent ground truth.

% ================================================================
\section{Problem formulation}
\label{sec:problem}

\subsection{Rough Heston and the Markovian lift}
Under rough Heston the spot price $S$ and its instantaneous variance $V$ satisfy
\begin{align}
  dS_t &= \sqrt{V_t}\,S_t\,dW^S_t, \\
  V_t  &= V_0 + \int_0^t K(t-s)\,\lambda(\theta - V_s)\,ds
          + \int_0^t K(t-s)\,\nu\sqrt{V_s}\,dW^V_s,
\end{align}
where $W^S, W^V$ are Brownian motions with correlation $\rho$, and
$K(t)=t^{\alpha-1}/\Gamma(\alpha)$ is the fractional kernel with
$\alpha=H+\tfrac12$ and $H<\tfrac12$ \citep{gatheral2018volatility,eleuch2019characteristic}.
The variance is a Volterra process: non-Markovian, non-semimartingale, and
admitting no finite-dimensional PDE.

The Markovian lift replaces $K$ with
$K(t)\approx\sum_{i=1}^{n}c_i e^{-x_i t}$, introducing $n$
Ornstein--Uhlenbeck factors $U^i$ with $dU^i_t = -x_i U^i_t\,dt + \sqrt{V_t}\,dW^V_t$
and $V_t = V_0 + \sum_i c_i U^i_t$
\citep{abijaber2019multifactor,abijaber2019lifting}.  The weights $(c_i,x_i)$
are placed on a geometric quadrature grid \citep{abijaber2019multifactor}; the
relative $L^2$ kernel error falls monotonically from $25.6\%$ at $n{=}5$ to
$0.39\%$ at $n{=}100$, so raising $n$ genuinely buys model fidelity.

\subsection{The lifted pricing PDE}
The discounted European call price $P(t,x,u)$, where $x=\log S$ and
$u=(u_1,\ldots,u_n)$ collects the factor values, satisfies the
$(n{+}1)$-dimensional backward Kolmogorov PDE
\begin{align}
  0 = {}& P_t + \!\left(r-\tfrac12 V\right)P_x + \tfrac12 V\,P_{xx}
         + \sum_{i=1}^{n}\!\left[-x_i u_i + \lambda(\theta-V)\right]P_{u_i}
         \nonumber\\
        & + \tfrac12\nu^2 V\sum_{i,j=1}^{n}P_{u_i u_j}
         + \rho\nu V\sum_{i=1}^{n} P_{x u_i} - rP ,
  \label{eq:pde}
\end{align}
with terminal condition $P(T,x,u)=\max(e^x - K, 0)$.
Because all factors share the single Brownian motion $W^V$, the full
double sum over $(i,j)$ and the $(n{+}1)$-dimensional cross term both appear;
the PDE dimension is exactly $n{+}1$.  At $n{=}4$ this is a five-dimensional
PDE; at $n{=}32$ it is thirty-three.

% ================================================================
\section{Network architecture and training}
\label{sec:arch}

\subsection{Baseline network}
We represent the price as a fully connected multilayer perceptron (MLP).  Every
hidden unit uses the $\tanh$ activation, which provides smooth, high-order
derivatives needed for automatic differentiation of the PDE residual; $\tanh$ is
standard for PINN residual training and avoids the dead-neuron pathology of ReLU
in the high-derivative setting \citep{raissi2019physics}.  All weight matrices are
initialised with Xavier normal initialisation \citep{glorot2010understanding}
and all biases with zeros.

The network maps
$(\tau, x, u_1,\ldots,u_n) \in \mathbb{R}^{2+n}$ to a scalar correction.  For the
single-point PINN (parameters fixed) we use width $W{=}64$ and depth $D{=}4$ hidden
layers; for the parametric PINN (five model parameters as additional inputs,
Section~\ref{sec:parametric}) we increase to $W{=}96$, $D{=}5$.
Table~\ref{tab:arch} summarises all architectural and training hyperparameters.

\begin{table}[t]
\centering
\caption{Architecture and training hyperparameters.  Values in parentheses are
for the parametric PINN (Section~\ref{sec:parametric}).}
\label{tab:arch}
\begin{tabular}{ll}
\toprule
\textbf{Component} & \textbf{Value} \\
\midrule
\multicolumn{2}{l}{\textit{Architecture}} \\
Network type & Fully connected MLP \\
Hidden layers & 4 (5 for parametric) \\
Neurons per layer & 64 (96 for parametric) \\
Activation & $\tanh$ (all hidden layers) \\
Weight initialisation & Xavier normal \\
Bias initialisation & Zero \\
Input dimension & $2 + n$ for single-point; $2 + n + 5$ for parametric \\
Output dimension & 1 (price correction $N_\theta$) \\
Floating-point precision & \texttt{float64} (PyTorch) \\
\midrule
\multicolumn{2}{l}{\textit{Optimisation}} \\
Optimiser & Adam \citep{kingma2015adam} \\
Initial learning rate & $10^{-3}$ \\
LR schedule & Cosine annealing to $0$ over all iterations \\
Gradient clipping & $\ell_2$ norm $\leq 1.0$ \\
Training iterations & 10,000--20,000 (single-point); 20,000--30,000 (parametric) \\
\midrule
\multicolumn{2}{l}{\textit{Collocation sampling (per iteration)}} \\
Interior collocation & 2,000 uniform + 500 near-strike (total 2,500) \\
Boundary points & 400 (split OTM / ITM) \\
Log-moneyness domain & $[\log K - 1.2,\;\log K + 1.2]$ \\
Factor sampling & MC-empirical mean / std (Section~\ref{sec:mcsamp}) \\
Tail coverage fraction & 20\% of $u$-points drawn at $2\times$ std \\
\midrule
\multicolumn{2}{l}{\textit{Loss weighting}} \\
Adaptive scheme & Learnable log-variance (homoscedastic uncertainty) \\
Terms balanced & PDE residual MSE, OTM boundary MSE, ITM boundary MSE \\
\midrule
\multicolumn{2}{l}{\textit{$\tau$-curriculum}} \\
Initial horizon & $\tau_0 = 0.15\,T$ \\
Ramp fraction & Linear growth to $T$ over first $50\%$ of iterations \\
\midrule
\multicolumn{2}{l}{\textit{Hardware}} \\
Platform & CPU (single machine); code is GPU-compatible (\texttt{torch.device}) \\
\bottomrule
\end{tabular}
\end{table}

\subsection{Hard-constrained initial-condition ansatz}
\label{sec:ansatz}
We train in time-to-maturity $\tau = T-t$ so the terminal condition becomes an
initial condition at $\tau{=}0$.  Rather than penalise it softly, we enforce it
exactly through the ansatz
\begin{equation}
  P(\tau,x,u) = P_{\mathrm{base}}(\tau,x) + \tau\,N_\theta(\tau,x,u),
  \label{eq:ansatz}
\end{equation}
where $N_\theta$ is the network output.  The $\tau$ prefactor forces the
learned correction to vanish at $\tau{=}0$ automatically, so the network never
fights the initial condition.  Without this, a soft penalty for the terminal
payoff oscillates against the PDE residual and prevents convergence.

\subsection{Integrated-variance baseline anchor}
\label{sec:baseline}
The baseline $P_{\mathrm{base}}$ is a Black--Scholes call price.  A critical
design choice is what volatility it is anchored at.  Anchoring at the
\emph{spot} variance $V_0$ leaves a systematic level bias: the factors mean-revert
from $V_0$ toward $\theta$ over the life of the option, so the model's actual
expected variance is higher than $V_0$.  We anchor instead at the model's expected
\emph{integrated} variance
\begin{equation}
  w(\tau) = \int_0^\tau \mathbb{E}[V_s]\,ds,
  \label{eq:intvar}
\end{equation}
computed by a short pre-flight Monte Carlo (the same simulation used for factor
collocation).  The baseline then uses $\sigma(\tau) = \sqrt{w(\tau)/\tau}$, which
correctly captures the mean-reversion effect; the network learns only the residual
stochastic-volatility correction.  This single change removed a uniform ${\sim}120$
vol\bp{} level bias we observed with the $V_0$ anchor.

\subsection{MC-informed factor collocation}
\label{sec:mcsamp}
The factor values $u$ at a given maturity $\tau$ are not centred at zero: they
mean-revert to a distribution set by the model dynamics.  Sampling from a fixed
symmetric box centred at zero causes the network to expend capacity on regions the
dynamics never visit, and---critically---to under-cover the region governing the
put-wing skew.  Before training we run a short lifted Monte Carlo simulation
(1,000 paths, 50 time steps) to estimate the mean and standard deviation of each
factor at the target maturity; training collocation is then drawn from
$\mathcal{N}(\hat{\mu}_i, \hat{\sigma}_i)$.  An additional 20\% of collocation
points are drawn at twice the estimated standard deviation to extend tail coverage.
This change cut smile RMSE by roughly $45\%$ (from ${\sim}210$ to ${\sim}114$ vol\bp)
relative to symmetric-box sampling.

\subsection{Self-adaptive loss weighting}
\label{sec:adaptive}
Three loss terms must be balanced at each training step: the PDE-residual MSE
$\mathcal{L}_{\mathrm{PDE}}$, and the MSEs at the deep-OTM ($P\to0$) and deep-ITM
($P\to S - Ke^{-r\tau}$) boundaries.  Fixed-weight balancing is brittle: the
ITM boundary target is $O(1)$ while the OTM target is $O(0)$, so a single lumped
weight over-emphasises the ITM term.  We adopt the learnable
homoscedastic-uncertainty scheme of \citet{kendall2018multitask} in the
self-adaptive spirit of \citet{mcclenny2023self}: each term $\mathcal{L}_k$ is
assigned a learnable log-variance $s_k$ and the objective is
\begin{equation}
  \mathcal{L} = \sum_{k\in\{\mathrm{PDE},\,\mathrm{OTM},\,\mathrm{ITM}\}}
               \bigl[e^{-s_k}\,\mathcal{L}_k + s_k\bigr].
  \label{eq:adaptive}
\end{equation}
The $s_k$ are optimised jointly with the network weights; large or noisy terms
are down-weighted automatically.  We never hand-tune a boundary weight.

\subsection{Causal \texorpdfstring{$\tau$}{tau}-curriculum}
\label{sec:curriculum}
PINNs that train simultaneously across all time-to-maturities can be pulled toward
spurious solutions at long $\tau$ before the short-$\tau$ region has converged,
because the residual at long $\tau$ depends causally on the short-$\tau$ solution
\citep{wang2024respecting}.  We address this with a causal curriculum: training
begins at $\tau_{\mathrm{max}} = 0.15\,T$ and linearly ramps to the full $T$ over
the first half of training.  The short-$\tau$ solution is therefore established
before the long-$\tau$ residual can perturb it.

\subsection{Training algorithm}
Algorithm~\ref{alg:train} summarises the full training procedure.

\begin{algorithm}[t]
\caption{PINN training for the lifted rough-Heston PDE}
\label{alg:train}
\begin{algorithmic}[1]
\Require Model parameters $(V_0,\lambda,\theta,\nu,\rho)$; lift $(c_i,x_i)_{i=1}^n$;
         maturity $T$; hyperparameters in Table~\ref{tab:arch}
\State Initialise MLP $N_\theta$ (Xavier normal); learnable $s = [s_1,s_2,s_3] \leftarrow 0$
\State \textbf{Pre-flight MC:} simulate lifted SDE, record $(\hat\mu_i, \hat\sigma_i)$
       for each factor and term structure $w(\tau)$ \Comment{Sections~\ref{sec:mcsamp},~\ref{sec:baseline}}
\For{iteration $t = 1, \ldots, N_{\mathrm{iter}}$}
  \State $\tau_{\mathrm{max}} \leftarrow \min\!\bigl(T,\; 0.15T + \tfrac{0.85T \cdot t}{0.5\,N_{\mathrm{iter}}}\bigr)$
         \Comment{$\tau$-curriculum, Sec.~\ref{sec:curriculum}}
  \State Sample interior: $\{(\tau_j, x_j, u_j)\}$ from
         $\mathrm{Unif}([0,\tau_{\mathrm{max}}])\times\mathrm{Unif}(x_{\mathrm{lo}},x_{\mathrm{hi}})
         \times\mathcal{N}(\hat\mu,\hat\sigma)$
  \State Sample boundary: $\{(\tau^{\mathrm{bc}}_j, x_{\mathrm{lo/hi}}, u^{\mathrm{bc}}_j)\}$
  \State $\hat P \leftarrow P_{\mathrm{base}}(\tau,x;\,w) + \tau\,N_\theta(\tau,x,u)$
         \Comment{Ansatz, Eq.~\eqref{eq:ansatz}}
  \State $\mathcal{L}_{\mathrm{PDE}} \leftarrow$ mean-squared PDE residual of $\hat P$
         (Eq.~\eqref{eq:pde})
  \State $\mathcal{L}_{\mathrm{OTM}} \leftarrow \|\hat P(\tau^{\mathrm{bc}},x_{\mathrm{lo}},u^{\mathrm{bc}})\|^2$,\quad
         $\mathcal{L}_{\mathrm{ITM}} \leftarrow \|\hat P(\tau^{\mathrm{bc}},x_{\mathrm{hi}},u^{\mathrm{bc}}) - (e^{x_{\mathrm{hi}}} - Ke^{-r\tau^{\mathrm{bc}}})\|^2$
  \State $\mathcal{L} \leftarrow \sum_k [e^{-s_k}\mathcal{L}_k + s_k]$
         \Comment{Adaptive weighting, Eq.~\eqref{eq:adaptive}}
  \State Backpropagate $\mathcal{L}$; clip gradients ($\ell_2 \leq 1.0$); Adam step
         on $(N_\theta, s)$; cosine LR anneal
\EndFor
\end{algorithmic}
\end{algorithm}

% ================================================================
\section{Evaluation methodology}
\label{sec:eval}

\subsection{Two independent ground truths}
We hold the network to two references validated independently.

\textbf{Fourier pricer.}  A semi-analytic pricer that solves the fractional Riccati
equation for the rough Heston characteristic function
\citep{eleuch2019characteristic} and inverts by Fourier integration
\citep{carr1999option}.  It never uses the lift, so it is the ground truth of the
\emph{true model} ($n\to\infty$); in the flat-volatility limit it reproduces
Black--Scholes to $3.7\times10^{-7}$ in price.

\textbf{Lifted Monte Carlo.}  An exponential-integrator Euler scheme that integrates
the $n$-factor lifted SDE with the $-x_i U^i$ drift handled exactly:
\begin{equation}
  U^i_{t+\Delta t} = e^{-x_i\Delta t}\,U^i_t
    + \int_t^{t+\Delta t}e^{-x_i(t+\Delta t - s)}\sqrt{V_s}\,dW^V_s.
  \label{eq:expint}
\end{equation}
The exponential-integrator is essential: the mean reversions $(x_i)$ span roughly
ten orders of magnitude for small $H$, and a naive Euler scheme diverges.  The
lifted MC matches Black--Scholes within Monte Carlo error (1,000 paths, 50 steps)
and agrees with the Fourier pricer for large $n$.  At finite $n$, the lifted MC is
the \emph{exact solution of the PDE the network solves}; the Fourier pricer is the
exact solution of the model.

\subsection{Error decomposition}
\label{sec:decomp}
Rather than reporting a single ``PINN-versus-Fourier'' figure, we decompose the
error per strike and in price space (where the signed contributions add exactly),
\begin{equation}
  \underbrace{(\mathrm{PINN}-\mathrm{Fourier})}_{\text{total}}
  =\underbrace{(\mathrm{PINN}-\mathrm{MC})}_{\text{solve-error}}
  +\underbrace{(\mathrm{MC}-\mathrm{Fourier})}_{\text{lift-error}}.
  \label{eq:split}
\end{equation}
The \emph{solve-error} measures how accurately the network solves the PDE it was
given; it is the quantity training controls.  The \emph{lift-error} is a model
approximation that only increasing $n$ can reduce; it is not the network's
responsibility.  Conflating the two gives a systematically pessimistic view of
network accuracy and leads to wrong engineering decisions (e.g.\ adding more
training iterations when the right action is to increase $n$).

We report price RMSE in basis points of spot (px\bp; $1\,\mathrm{px\,bp} =
10^{-4}\times S$) where \eqref{eq:split} is additive.  Implied-vol equivalents
(vol\bp) are quoted where they aid interpretation.  Note that the two components
carry opposite signs, so their RMSE values do not add.

\subsection{Homogeneity re-scaling}
We train the network at fixed strike $K{=}1$ and read off the smile $C(1,K)$
through the call's degree-one homogeneity:
\begin{equation}
  C(1,K) = K\,C\!\left(1/K,\;1\right),
  \label{eq:homog}
\end{equation}
which holds exactly because the variance dynamics are independent of the spot
level.  Dropping this re-scaling and comparing mismatched options inflates the
apparent error to ${\sim}2700$ vol\bp; we verified this is a comparison artefact,
not a modelling failure.

% ================================================================
\section{Calibration procedures and data}
\label{sec:calib}

\subsection{Data}
We use end-of-day BTC option quotes from Deribit, a cryptocurrency derivatives
exchange and data provider, retrieved through its public API (no API key required).
After filtering for positive bid--ask spreads and non-zero open interest, the
surface covers 13 maturities and 97 points spanning a moneyness range of roughly
$0.6$--$1.7$.  We set $r{=}0$ (standard for crypto, since there is no clean carry)
and infer the forward from the put--call parity of liquid near-the-money quotes.
Implied volatilities are computed from mid prices with a robust Newton--Raphson
Black--Scholes inverter.  The calibration objective throughout is the
vega-weighted implied-volatility RMSE.

\subsection{Calibration methods}
We compare three calibration routes head to head.
\begin{enumerate}
  \item \textbf{Fourier-in-the-loop.}  Differential evolution (60 generations,
        all 5 parameters) with the Fourier pricer inside the objective, followed
        by Nelder--Mead polish.  This is our from-scratch reference.
  \item \textbf{PINN-only.}  Gradient descent on the model parameters using the
        parametric network as a differentiable pricer, with zero Fourier calls.
  \item \textbf{Hybrid.}  The parametric PINN's single-shot solution provides the
        initial guess; a short Nelder--Mead Fourier polish then sharpens it.
\end{enumerate}

% ================================================================
\section{Results}
\label{sec:results}

\subsection{The network scales where the grid cannot}
\label{sec:nconv}
Our headline experiment trains one architecture (width 64, depth 4) at
$n{=}4,8,16,32$ on the canonical BTC calibration at $T{=}0.15$, over three
independent seeds (Table~\ref{tab:nconv}, Figure~\ref{fig:nconv}).  The solve-error
against the lifted MC grows only mildly: from $5.4\pm1.0$ to $10.7\pm2.0$ px\bp---a
doubling over an eight-fold rise in dimension.  We resist calling it flat: an early
single-seed run gave that impression, and three-seed statistics corrected us.  With
three seeds the error bars at adjacent $n$ overlap, so the robust statement is the
two-fold growth across the full range.  Wall-clock cost stays bounded throughout
(6--10 minutes on CPU, per-iteration time essentially flat at 55--62\,ms) because
the factor-Hessian computation is vectorised over the full double sum.

Compare this against a finite-difference alternative.  A naive full-tensor grid on
the same $(n{+}1)$-dimensional PDE at a modest 50 nodes per dimension needs
$50^{n+1}$ points: $3.1\times10^{8}$ at $n{=}4$, $2.0\times10^{15}$ at $n{=}8$,
and $1.2\times10^{56}$ at $n{=}32$.  Even $n{=}8$---roughly $16$ petabytes in double
precision---is out of reach, while the network trains every $n$ on a laptop.
Structured schemes (sparse grids, ADI, low-rank tensor-train) reach further but
still degrade with $n$; none is standard or validated for this coupled
$\rho$--$\nu$ PDE.  The conservative engineering verdict is that a competitive
grid solve at $n{=}16$--$32$ is \emph{impractical}, not provably impossible.

\begin{table}[t]
\centering
\caption{Dimension scaling: single-point PINN at $T{=}0.15$ on the canonical BTC
calibration.  Solve-error is 3-seed mean\,$\pm$\,SD in price basis points of spot;
lift-error and FD grid size are deterministic.}
\label{tab:nconv}
\begin{tabular}{r r r r r}
\toprule
$n$ & solve-error (px\bp) & lift-error (px\bp) & train time (s) & FD grid $50^{(n+1)}$\\
\midrule
4  & $5.4\pm1.0$  & $16.4$ & $397$ & $3.1\times10^{8}$  \\
8  & $6.1\pm1.4$  & $7.9$  & $386$ & $2.0\times10^{15}$ \\
16 & $8.8\pm1.6$  & $3.5$  & $591$ & $7.6\times10^{28}$ \\
32 & $10.7\pm2.0$ & $1.7$  & $457$ & $1.2\times10^{56}$ \\
\bottomrule
\end{tabular}
\end{table}

\begin{figure}[t]\centering
\includegraphics[width=0.82\linewidth]{fig1_n_convergence.png}
\caption{Dimension scaling.  PINN solve-error (blue, 3-seed mean\,$\pm$\,SD)
grows sub-linearly and stays under ${\sim}11$ px\bp{} to $n{=}32$, while the
lift-error (red) falls monotonically and training time (green) stays bounded.
The inset shows the naive finite-difference grid size $50^{n+1}$, which is
infeasible by $n{=}8$.}
\label{fig:nconv}
\end{figure}

\subsection{Error decomposition and a design rule for $n$}
\label{sec:decomp_results}
The decomposition \eqref{eq:split} overturned how we read our own results.  A
single ``PINN vs Fourier'' implied-vol RMSE of $114$--$127$ vol\bp{} at $n{=}4$
\emph{looks} like a network failure.  Split by \eqref{eq:split}, almost all of
it is the lift: $37$ vol\bp{} of solve-error on top of $133$ of lift (they partly
cancel to a total of $121$); with $\rho{=}0$ the solve-error is only $8$ vol\bp{}
while the lift contributes the same $133$.  The network solves its PDE to well
under $50$ vol\bp{}; what had appeared to be a modelling failure of the network
was the finite-$n$ model catching up to reality.

Because the lift-error shrinks monotonically with $n$ ($133, 63, 27, 12$ vol\bp{}
at $n{=}4,8,16,32$) while the solve-error grows gently, the two curves cross.
At $T{\approx}0.15$ the crossover falls near $n{=}8$--$10$
(Figure~\ref{fig:nconv}), and this gives a concrete engineering design rule:
\emph{below the crossover, add factors; above it, stop}.  The crossover shifts
with maturity: the lift-error grows with $T$ (from $8.6$ to $31.8$ px\bp{} at
$n{=}4$ as $T$ rises from $0.05$ to $0.5$), so longer-maturity pricing requires
larger $n$.

Robustness sweeps over $T\in\{0.05,0.15,0.5\}$ and $H\in\{0.05,0.10,0.15\}$
confirm that the solve-error stays small and bounded everywhere ($1$--$9.5$ px\bp{}
across all configurations), while the lift-error varies predictably.  The network
does carry a mild directional bias at large $n$: at $n{=}32$ it systematically
over-prices ATM and call strikes by roughly $+87$ to $+114$ vol\bp{} while the
put wing is near zero; this bias grows with $n$ and is the main watch item for
future high-$n$ deployments.

\subsection{Bitcoin is genuinely rough}
\label{sec:btc}
Calibrated to the full Deribit BTC surface with a full-budget optimiser
(differential evolution, 60 generations, all 13 maturities, 97 points, followed
by Nelder--Mead polish), rough Heston returns
\begin{equation*}
  H{=}0.090,\quad \rho{=}{-}0.792,\quad V_0{=}0.103,\quad
  \theta{=}0.253,\quad \lambda{=}2.567,\quad \nu{=}0.300,
\end{equation*}
at a vega-weighted implied-vol RMSE of 496 vol\bp{} (Figure~\ref{fig:calib}).
The Hurst exponent is firmly rough and the correlation strongly negative---the
signature of a crypto market with a steep, persistent left skew.  We do not dress
up the fit quality: at 496 vol\bp{} it is coarse, and the per-maturity smiles in
Figure~\ref{fig:calib}(b) show where the model strains.  That coarseness is a
statement about the model's expressiveness on wide, richly curved crypto smiles,
not about our solver; it is orthogonal to the solve-error studied above.

\begin{figure}[t]\centering
\includegraphics[width=0.98\linewidth]{fig2_calibration_fit.png}
\caption{Rough Heston calibrated to the Deribit BTC surface (97 points, 13
maturities).  (a)~Market versus model implied volatility at all points (dashed:
identity; colour: maturity).  (b)~Per-maturity smiles for four representative
maturities: market (markers) versus model (dashed).  The residual spread reflects
a model-expressiveness limitation on wide crypto smiles, not a solver artefact.}
\label{fig:calib}
\end{figure}

\subsection{Parametric PINN and the warm-start hybrid}
\label{sec:parametric}
A single network cannot cover both the breadth of a parameter box and the accuracy
of a point-trained model.  Table~\ref{tab:calib} and Figure~\ref{fig:hybrid}
report the calibration experiment with three seeds.

The parametric network adds $(H,\nu,\rho,\lambda,\theta)$ as five additional inputs,
holding $V_0$ and the maturity fixed; we use $n{=}10$ (the crossover rule,
Section~\ref{sec:decomp_results}) and width 96, depth 5.  At the box centre its
solve-error is $11.0\pm4.7$ px\bp{} versus $5.9\pm1.2$ for a single-point network
at the same $n$---roughly a two-fold penalty.  The penalty is uneven: it hovers near
$9$--$11$ px\bp{} at typical parameter vectors but climbs to $47.7$ in the rough
low-$H$ corner and $29.8$ at high $\nu$, where the PDE is stiffest.

Used as a standalone calibrator the parametric PINN is fast (approximately $2$\,s,
zero Fourier calls) but poor: $222\pm20$ vol\bp{} under the true model.  Two
factors explain this.  First, the network fits its own imperfect prices, not the
market.  Second, the single-slice BTC optimum prefers $\nu\approx0.87$ and
$\theta\approx0.39$, which lie \emph{outside} the training box, so the network
cannot reach the true minimum.  We would not recommend the standalone parametric
PINN as a calibrator.

Used as a warm start it is effective.  Seeding the Nelder--Mead Fourier polish
with the PINN's parameter guess reaches $29.4\pm7.1$ vol\bp{}---on par with the
from-scratch result on a 34--52\% smile---in only $237\pm46$ Fourier evaluations
and ${\approx}23$\,s: roughly nine times fewer true-model calls and seven times
faster than from scratch.  The value proposition is amortised: train once,
then recalibrate cheaply, with the exact pricer kept in the loop for a brief,
decisive finish.

\begin{table}[t]
\centering
\caption{Calibration of the BTC $T{\approx}0.15$ slice, three methods, 3 PINN seeds.
All fit-quality figures are judged by the true Fourier model over the same feasible
region.}
\label{tab:calib}
\begin{tabular}{l c c c}
\toprule
Method & Fit RMSE (vol\bp) & Wall-clock (s) & Fourier evals \\
\midrule
Fourier-loop from scratch         & $23.7$        & $168$   & $2018$         \\
PINN-only (single-shot)           & $222\pm20$    & $\sim2$ & $0$            \\
Hybrid (PINN + Fourier polish)    & $29.4\pm7.1$  & $23$    & $237\pm46$     \\
\bottomrule
\end{tabular}
\end{table}

\begin{figure}[t]\centering
\includegraphics[width=0.82\linewidth]{fig3_hybrid.png}
\caption{Hybrid calibration.  Fit quality (left) and Fourier-evaluation count
(right) for the three methods.  The hybrid matches from-scratch quality at
${\sim}9\times$ fewer Fourier evaluations; the PINN alone is fast but inaccurate.}
\label{fig:hybrid}
\end{figure}

% ================================================================
\section{Discussion}
\label{sec:disc}

\paragraph{What the scaling result does and does not say.}
The finding we care about most is that the network's own error in solving the
lifted PDE degrades only gently as the dimension grows, while grid methods become
impractical.  We state it in its honest, weak form: the solve-error is not
constant in $n$---a claim an early single-seed run seemed to support and three-seed
statistics refuted---it grows sub-linearly, roughly doubling from $n{=}4$ to
$n{=}32$.  And the grid comparison uses the naive $50^{n+1}$ count; sparse-grid,
ADI, and tensor-train methods go further, so the defensible claim is
\emph{impractical, not impossible}.  Even so, the contrast---sub-linear error and
bounded wall-clock versus a grid that cannot be stored---is a concrete instance of
the mesh-free advantage, and we believe it is the clearest such instance on an
economically real, high-dimensional PDE.

\paragraph{Four design choices are load-bearing.}
We ablated each stabiliser explicitly.  Without the hard-constrained ansatz, the
soft IC penalty oscillates against the residual and convergence stalls.  Without the
integrated-variance anchor, a uniform $120$ vol\bp{} level bias remains.  Without
MC-informed factor sampling, the smile flattens and RMSE roughly doubles.  Without
self-adaptive weighting or the $\tau$-curriculum, seed variance is markedly higher
(from $339\pm62$ to $310\pm28$ vol\bp{} in an early stability comparison).  Each
of these design choices is transferable to other high-dimensional PDE-pricing
problems.

\paragraph{A lesson about validation methodology.}
The error decomposition revealed something methodologically important.  When a
surrogate replaces an exact solver, it is dangerously easy to validate it against
the wrong reference.  Our network solves the finite-$n$ lifted PDE, whose exact
solution is the lifted Monte Carlo, not the $n\to\infty$ model; comparing straight
to the true model conflates solve-error with lift-error.  We found the two carry
opposite signs and that the historical ``PINN-versus-Fourier'' gap was
lift-dominated, not the network's failure.  \emph{Validate against a
same-fidelity reference, report the two errors separately, and let their crossover
set the approximation order}: none of this is special to rough volatility, and we
urge it on anyone placing a learned solver on top of an approximated operator.

\paragraph{Engineering trade-offs of the parametric PINN.}
One network spread over five parameters is markedly less accurate than a
point-trained one, worst in the rough, high-vol-of-vol corners, and calibrating
with it alone is fast but unreliable.  Yet the same network is a fine warm start.
The value proposition is amortised: the offline training cost is paid once, and
each subsequent recalibration is roughly an order of magnitude cheaper.  A larger
training budget would narrow the corner penalty.  The key limitation today is that
the training box must contain the market optimum; box design therefore requires a
prior characterisation of the market regime.

\paragraph{Limitations and watch items.}
Our evidence rests on a single market snapshot (one Deribit BTC surface), and the
parametric and calibration studies fix the maturity to a single slice.  Timings
are single-machine CPU figures at moderate budgets.  The mild directional
solve-bias at large $n$ (ATM and call over-pricing growing to ${\sim}100$ vol\bp{}
at $n{=}32$) is a watch item for high-$n$ deployments.  The sub-linear trend need
not continue arbitrarily far.

\paragraph{Future work.}
The natural extensions are: (i)~add the maturity as a network input for whole-surface
pricing; (ii)~scale width, depth, and GPU training to shrink the parametric corner
penalty and probe larger $n$; (iii)~extend to other rough models (rough Bergomi)
and to American and path-dependent payoffs; (iv)~investigate whether residual
reweighting or a correction term resolves the high-$n$ directional bias;
(v)~benchmark the mesh-free solver directly against structured finite differences
at moderate $n$ where the latter still stand a chance.

% ================================================================
\section{Conclusion}
\label{sec:conc}

We have presented a physics-informed neural network for the $(n{+}1)$-dimensional
Markovian-lifted rough-Heston pricing PDE, with a complete description of the
architecture and four training stabilisers that are each necessary for stable
convergence.  The network's solve-error grows only sub-linearly in the lifted
dimension---staying below eleven price basis points to $n{=}32$ on the BTC
calibration---while naive finite-difference grids become infeasible by $n{=}8$.
Separating solve-error from lift-error is essential to read the numbers correctly
and yields a concrete, maturity-dependent crossover rule for the factor count.
Applied to live Bitcoin options the model is genuinely rough ($H{\approx}0.09$), and
a parametric PINN, while not accurate enough to calibrate independently, is an
effective warm start that recalibrates at roughly nine times fewer true-model
evaluations.  The engineering message is that where grid methods succumb to
dimension, a mesh-free PINN degrades gracefully; its greatest practical value is
as a differentiable warm start for fast, high-quality recalibration.

% ================================================================
\section*{CRediT authorship contribution statement}

\textbf{Aviral Sharma:} Conceptualisation, Methodology, Software, Validation,
Formal analysis, Writing -- original draft, Writing -- review \& editing.
\textbf{Kamal Preet Singh:} Writing -- review \& editing, Funding acquisition.
\textbf{Vipul Sharma:} Writing -- review \& editing.

\section*{Declaration of competing interests}

The authors declare that they have no known competing financial interests or
personal relationships that could have appeared to influence the work reported in
this paper.

\section*{Data availability}

The BTC option data were retrieved from the Deribit public REST API
(no API key required) at \texttt{https://www.deribit.com/api/v2}.  The raw data
snapshot used in the experiments, the calibrated parameter file, and the complete
PyTorch implementation are available at \todo{repository URL upon acceptance}.

\section*{Declaration of generative AI and AI-assisted technologies in the
writing process}

During the preparation of this manuscript the authors used Claude (Anthropic)
to assist with prose drafting and structural organisation.  After using this
tool, the authors reviewed and edited all content as needed and take full
responsibility for the accuracy and integrity of the published work.  The
mathematical derivations, experimental results, code, and all scientific claims
are the authors' own.

\section*{Acknowledgements}

This research received no specific grant from any funding agency in the public,
commercial, or not-for-profit sectors.

% ================================================================
%% The .bbl file is pre-generated (eaai_v3.bbl) so the paper compiles in a
%% single pdflatex pass without running BibTeX.
\bibliographystyle{elsarticle-harv}
\bibliography{refs}

\end{document}
